\documentclass[../../main]{subfiles}

\ifSubfilesClassLoaded{%
    \zexternaldocument*{../1.Category-Theory/exported/aux/main}%
    \zexternaldocument*{../2.Homological-Algebra/exported/aux/main}%
    \zexternaldocument*{../3.Sheaves/exported/aux/main}%
    \zexternaldocument*{1.Cohomology-of-sheaves}%
    \zexternaldocument*{2.Cohomology-and-homotopy}%
}{}

\begin{document}
\chapter{Cohomology theories}\label{ch:cohomology-theories}


Historically, Eilenberg and Steenrod showed that the various forms of cohomology could be characterized uniquely by a small set of simple axioms.
The goal of this chapter is to verify that that \hyperref[def:sheaf-cohomology]{sheaf cohomology} satisfies these axioms.

In this chapter we restrict ourselves to \hyperref[def:sheaf]{sheaves of abelian groups}, i.e: we take \(\mathbf{k} = \mathbb{Z}\).

\section{Cohomology theories and CW-Complexes}

\begin{definition}\label{def:topological-pair}
  A \textbf{topological pair} is a tuple \((X, A)\) where \(X\) is a topological space and \(A\) is a subspace of \(X\).
  A morphism of topological pairs \((X, A) \to (Y, B)\) is a continuous map \(f\colon X \to Y\) such that \(f(A) \subseteq B\).
  
  We denote by \(\Top_{2}\) the \hyperref[def:category]{category} of topological pairs.
\end{definition}

We start by giving the definition of a cohomology theory following~\cite{eilenbergAxiomaticApproachHomology1945}.
\begin{definition}\label{def:cohomology-theory}
  A \textbf{cohomology theory} (on \(\Top\)) consists of
  \begin{enumerate}
    \item A sequence of \hyperref[def:functor]{contravariant functors} \(\cohomology{n}\colon \Top_{2} \to \Ab\) for all \(n \in \mathbb{Z}\).
          
          We denote by \(\cohomology{n}(X) \coloneq \cohomology{n}(X, \varnothing)\).
          If \(f\colon (X, A) \to (Y, B)\) is a continuous map we denote \(f^{\ast}\colon \cohomology{n}(Y, B) \to \cohomology{n}(X, A)\) the induced map (where \(n\) is to be understood from the context).
    \item A sequence of \hyperref[def:natural-transformation]{natural transformations} \(\delta^{n}\colon \cohomology{n}(A) \to \cohomology{n+1}(X, A)\).
  \end{enumerate}
  Satisfying the following axioms:
  \begin{enumerate}
    \item \textbf{Homotopy}: If \(f_{0}, f_{1}\colon (X, A) \to (Y, B)\) are homotopic then its induced maps on \(\cohomology{n}\) are equal
          \begin{equation*}
            f_{0}^{\ast} = f_{1}^{\ast}\colon \cohomology{n}(Y, B) \to \cohomology{n}(X, A)
          \end{equation*}
    \item \textbf{Excission}: If \((X, A) \in \Top_{2}\) and \(U \subseteq A\) is such that \(\operatorname{ad}(U) \subset \operatorname{int}(A)\) then the inclussion map \(i\colon (X \setminus U, A \setminus U) \to (X, A)\) induces an isomorphism in \(\cohomology{n}\):
          \begin{equation*}
            i^{\ast}\colon \cohomology{n}(X, A) \isoarrow \cohomology{n}(X \setminus U, A \setminus U)
          \end{equation*}
    \item \textbf{Exactness}: Each pair \((X, A) \in \Top_{2}\) induces a \hyperref[def:exact-sequence]{long exact sequence} via the inclussions \(i\colon A \to X\) and \(j\colon X \to (X, A)\):
          \begin{center}
            \begin{tikzcd}
              \cdots \arrow[r] & \cohomology{n-1}(X) \arrow[r, "i^{\ast}"] \arrow[d, phantom, ""{coordinate, name=X}] & \cohomology{n-1}(A) \arrow[dll, rounded corners, "\delta^{n-1}" description, to path={
                -- ([xshift=2ex]\tikztostart.east)
                |- (X) [near end]\tikztonodes
                -| ([xshift=-2ex]\tikztotarget.west)
                -- (\tikztotarget)}] \\
              \cohomology{n}(X, A) \arrow[r, "j^{\ast}"] & \cohomology{n}(X) \arrow[r, "i^{\ast}"] \arrow[d, phantom, ""{coordinate, name=Y}] & \cohomology{n}(A) \arrow[dll, rounded corners, "\delta^{n}" description, to path={
                -- ([xshift=2ex]\tikztostart.east)
                |- (Y) [near end]\tikztonodes
                -| ([xshift=-2ex]\tikztotarget.west)
                -- (\tikztotarget)}] \\
              \cohomology{n+1}(X, A) \arrow[r, "j^{\ast}"] & \cohomology{n+1}(X) \arrow[r] & \cdots
            \end{tikzcd}
          \end{center}
  \end{enumerate}
  We also require the following axiom due to Milnor~\cite{mllnorAXIOMATICHOMOLOGYTHEORY}
  \begin{enumerate}
    \item[4.] \textbf{Additivity}: If \(X = \coprod_{i \in \mathcal{I}} X_{i}\) is the disjoint union of topological spaces then
          \begin{equation*}
            \cohomology{n}(X) \cong \prod_{i \in \mathcal{I}} \cohomology{n}(X_{i})
          \end{equation*}
  \end{enumerate}

  We say that a cohomology theory is ordinary if it further verifies
  \begin{enumerate}
    \item[5.] \textbf{Dimension}: If \(\{\ast\}\) is the one-point topological space then for all \(n \neq 0\)
          \begin{equation*}
            \cohomology{n}(\{\ast\}) = 0
          \end{equation*}
          In this case, \(\cohomology{0}(\{\ast\})\) is called the coefficient group of the theory.
  \end{enumerate}
\end{definition}

For sufficiently nice spaces, one can show that there exists only one \hyperref[def:cohomology-theory]{ordinary cohomology theory} with a given coefficient group.

\begin{definition}[CW-complex]\label{def:CW-complex}
  A \textbf{CW-complex} is a topological space \(X\) constructed inductively in the following way:
  \begin{enumerate}
    \item Start with a discrete set \(X^{0}\), called the 0-cells of \(X\).
    \item Inductively, form the \(n\)-skeleton \(X^{n}\) from \(X^{n-1}\) by attaching \(n\)-cells \(e_{\alpha}^{n}\) via maps \(\varphi_{\alpha}\colon \mathbb{S}^{n-1} \to X^{n-1}\).

          This means that \(X^{n}\) is the quotient space of \(X^{n-1} \coprod_{\alpha} \mathbb{D}_{\alpha}^{n}\) under the identifications \(x \sim \varphi_{\alpha}(x)\) for \(x \in \partial \mathbb{D}_{\alpha}^{n-1}\).
          The cell \(e_{\alpha}^{n}\) is the homeomorphic image of \(\mathbb{D}_{\alpha}^{n} \setminus \partial \mathbb{D}_{\alpha}^{n-1}\) under the quotient map.
    \item Set \(X \coloneq \bigcup_{n \in \mathbb{N}} X^{n}\) with the weak topology: a set \(A \subseteq X\) is open if \(A \cap X^{n}\) is open in \(X^{n}\) for all \(n\).
  \end{enumerate}
\end{definition}

We will not prove any properties of \hyperref[def:CW-complex]{CW-complexes}, refering the reader to \cite[Appendix]{hatcherAlgebraicTopology2002}.
A particularly important result is that any compact manifold can be homotopied into a \hyperref[def:CW-complex]{CW-complex}.
\begin{proposition}[{\cite[Corollary A.12]{hatcherAlgebraicTopology2002}}]\label{prop:manifold-is-CW-complex}
  A compact manifold is homotopy equivalent to a \hyperref[def:CW-complex]{CW-complex}.
\end{proposition}

\begin{theorem}[{Uniqueness of cohomology theories~---~\cite[Theorem 4.59]{hatcherAlgebraicTopology2002}}]\label{thm:uniqueness-cohomology-theories}
  Let \(\{\cohomology{n}, \delta^{n}\}_{n \in \mathbb{N}}\) and \(\left\{ \widetilde{\operatorname{H}}^{n}, \widetilde{\delta}^{n} \right\}_{n \in \mathbb{N}}\) be two \hyperref[def:cohomology-theory]{ordinary cohomology theories} in the category of \hyperref[def:CW-complex]{CW pairs} with the same coefficient group.
  Then for all \(n \geq 0\) there is a \hyperref[def:natural-transformation]{natural isomorphism}
  \begin{equation*}
    \cohomology{n}(X, A) \cong \widetilde{\operatorname{H}}^{n}(X, A)
  \end{equation*}
\end{theorem}

The most famous \hyperref[def:cohomology-theory]{cohomology theory} is \textbf{singular cohomology} \(\operatorname{H}_{\operatorname{sing}}\).
In fact the axioms of \hyperref[def:cohomology-theory]{cohomology theories} were modelled after its behaviour and the proof of \autoref{thm:uniqueness-cohomology-theories} works by proving that any \hyperref[def:cohomology-theory]{ordinary cohomology theory} is isomorphic to singular cohomology (given that they have the same coefficient group).

We will not go into any details of its definition or properties, refering the reader to \cite[Chapter 3]{hatcherAlgebraicTopology2002} for more details.

\begin{remark}
  Another strategy to prove that \hyperref[def:sheaf-cohomology]{sheaf cohomology} is an \hyperref[def:cohomology-theory]{ordinary cohomology theory} is to prove that it coincides with singular cohomology (or any other ordinary cohomology theory) by finding a suitable resolution as we did for the de Rham case.

  This however would require knowledge about the chosen \hyperref[def:cohomology-theory]{cohomology theory} (and of course proving that it is indeed an ordinary cohomology theory).
  For the interested reader we refer to \cite[Chapter III]{bredonSheafTheory1997}.
\end{remark}


\section{Relative sheaf cohomology}

The definition of a \hyperref[def:cohomology-theory]{cohomology theory} is done in the \hyperref[def:topological-pair]{category of topological pairs} \(\Top_{2}\).
However, our definition of \hyperref[def:sheaf-cohomology]{sheaf cohomology} only accepts a single space as an input.

To fix this, we introduce the following version of cohomology.

\begin{definition}[Relative Sheaf Cohomology]\label{def:relative-sheaf-cohomology}
  Let \(X\) be a topological space and \(A \subseteq X\) a closed subset.
  Let \(j\colon X \setminus A \hookrightarrow X\) be the open immersion of the complement.
  For any \hyperref[def:sheaf]{sheaf of \(\mathbf{k}\)-modules} \(\sheaf{F}\) over \(X\), the \textbf{relative cohomology} of the pair \((X, A)\) with coefficients in \(\sheaf{F}\) is:
  \begin{equation*}
    \cohomology{i}(X, A; \sheaf{F}) \coloneq \cohomology{i}(X; j_{!}j^{-1}\sheaf{F})
  \end{equation*}

  If \(A = \varnothing\) we recover the usual \hyperref[def:sheaf-cohomology]{sheaf cohomology}
  \begin{equation*}
    \cohomology{i}(X, \varnothing; \sheaf{F}) = \cohomology{i}(X; \sheaf{F})
  \end{equation*}
\end{definition}

\begin{remark}
  \hyperref[def:relative-sheaf-cohomology]{Relative sheaf cohomology} is the \hyperref[def:right-derived-functor]{right derived functor}
  \begin{equation*}
    \Gamma(X; j_{!}j^{-1}\sheaf{F}) \coloneq \left\{ s \in \Gamma(X, \sheaf{F}) \mid s|_{A} = 0\right\}
  \end{equation*}
\end{remark}

Now we are ready to prove that \hyperref[def:relative-sheaf-cohomology]{relative sheaf cohomology} is a \hyperref[def:cohomology-theory]{cohomology theory} by verifying the axioms one by one.

Note that the dimension axiom was already verified back in \autoref{ex:cohomology-over-point}.

\begin{theorem}[Homotopy]\label{thm:homotopy-pairs}
  Let \(f_{0}, f_{1}\colon (X, A) \rightrightarrows (Y, B)\) be two \hyperref[def:homotopy-over]{homotopic maps of pairs}.
  Let \(\underline{M}_{Y} \in \Mod( \underline{\mathbf{k}}_{Y} )\) be a \hyperref[def:constant-sheaf]{constant sheaf} of \(\mathbf{k}\)-modules over \(Y\).
  Then for all \(i \geq 0\), the induced maps on \hyperref[def:relative-sheaf-cohomology]{relative cohomology} coincide:
  \begin{equation*}
    f_{0}^{\ast} = f_{1}^{\ast} \colon \cohomology{i}(Y, B; \underline{M}_{Y}) \to \cohomology{i}(X, A; \underline{M}_{X})
  \end{equation*}
\end{theorem}
\begin{proof}
  Let \(H\colon (X \times I, A \times I) \to (Y \times I, B \times I)\) be the homotopy between \(f_{0}\) and \(f_{1}\).
  We have
  \begin{equation*}
    f_{0}^{\ast} = j_{0}^{\ast} \circ H \qquad f_{1}^{\ast} = j_{1}^{\ast} \circ H
  \end{equation*}
  Therefore we need to prove
  \begin{equation*}
    j_{0}^{\ast} = j_{1}^{\ast}\colon \cohomology{i} \left( X \times I, A \times I; \underline{M}_{X \times I} \right) \to \cohomology{i} \left( X, A; \underline{M}_{X} \right)
  \end{equation*}
  Note that \(p_{(X, A)} \circ j_{0} = \id_{(X, A)} = p_{(X, A)} \circ j_{1}\).
  Furthermore
  \begin{equation*}
    i_{0} \circ p_{(X, A)}\colon (X \times I, A \times I) \to (X \times I, A \times I)
  \end{equation*}
  is homotopic to the identity.

  Therefore by \autoref{thm:homotopy-invariance-cohomology} we get isomorphisms
  \begin{align*}
    p_{X}^{\ast}\colon& \cohomology{i}(X; \underline{M}_{X}) \isoarrow \cohomology{i} \left( X \times I; \underline{M}_{X \times I} \right) \\
    p_{A}^{\ast}\colon& \cohomology{i}(A; \underline{M}_{A}) \isoarrow \cohomology{i} \left( A \times I; \underline{M}_{A \times I} \right)
  \end{align*}
  By the \hyperref[thm:exactness-axiom]{exactness axiom} we have \hyperref[def:exact-sequence]{long exact sequences}
  \begin{center}
    \begin{tikzcd}[column sep=small]
      \cdots \arrow[r] & \cohomology{i-1}(A; \underline{M}) \arrow[r] \arrow[d, "p_{A}^{\ast}"', "\sim"] & \cohomology{i}(X, A; \underline{M}) \arrow[r] & \cohomology{i}(X; \underline{M}) \arrow[r] \arrow[d, "p_{X}^{\ast}"', "\sim"]& \cohomology{i}(A; \underline{M}) \arrow[r] \arrow[d, "p_{A}^{\ast}"', "\sim"]& \cdots \\
      \cdots \arrow[r] & \cohomology{i-1}(A \times I; \underline{M}) \arrow[r] & \cohomology{i}(X \times I, A \times I; \underline{M}) \arrow[r] & \cohomology{i}(X \times I; \underline{M}) \arrow[r] & \cohomology{i}(A \times I; \underline{M}) \arrow[r] & \cdots
    \end{tikzcd}
  \end{center}
  Therefore by the \hyperref[lemma:five-lemma]{five lemma}
  \begin{equation*}
    p_{(X, A)}^{\ast}\colon \cohomology{i}(X, A; \underline{M}) \isoarrow \cohomology{i}(X \times I, A \times I; \underline{M})
  \end{equation*}
  By functoriallity of cohomology we have
  \begin{equation*}
    j_{0}^{\ast} \circ p_{(X, A)}^{\ast} = \id = j_{1}^{\ast} \circ p_{(X, A)}^{\ast}
  \end{equation*}
  Therefore since \(p_{(X, A)}^{\ast}\) is an isomorphism
  \begin{equation*}
    j_{0}^{\ast} = j_{1}^{\ast} \qedhere
  \end{equation*}
\end{proof}

\begin{theorem}[Excision]\label{thm:excision-axiom}
  Let \(A \subseteq X\) be a closed subset.
  Let \(Z \subseteq A\) be a subset such that \(\operatorname{ad}(Z) \subseteq \operatorname{int}(A)\).
  Then for any \hyperref[def:sheaf]{sheaf} \(\sheaf{F} \in \Mod(\underline{\mathbf{k}}_{X})\):
  \begin{equation*}
    \cohomology{i}(X, A; \sheaf{F}) \cong \cohomology{i}(X \setminus Z, A \setminus Z; \sheaf{F})
  \end{equation*}
\end{theorem}
\begin{proof}
  Let \(U \coloneq X \setminus A\) and \(j\colon U \hookrightarrow X\) the open immersion.
  By \hyperref[def:relative-sheaf-cohomology]{definition} we have
  \begin{equation*}
    \cohomology{i}(X, A; \sheaf{F}) = \cohomology{i}(X; \sheaf{G})
  \end{equation*}
  where \(\sheaf{G} \coloneq j_{!}j^{-1}\sheaf{F}\).

  Without loss of generality suppose that \(Z\) is closed (if not consider \(\operatorname{ad}(Z)\)).
  Let \(V = X \setminus Z\) and \(W = \operatorname{int}(A)\).
  Since \(Z\) is closed \(V\) is open. Furthermore since \(Z \subseteq \operatorname{int}(A) = W\) we have an open cover \(X = V \cup W\).

  By \hyperref[thm:mayer-vietoris-open]{Mayer-Vietoris} we have a \hyperref[def:exact-sequence]{long exact sequence}
  \begin{center}
    \begin{tikzcd}
      {} & {} & \cdots \arrow[r] \arrow[d, phantom, ""{coordinate, name=X}] & \cohomology{i-1}(V \cap W; \sheaf{G}|_{V \cap W}) \arrow[dll, rounded corners, to path={
        -- ([xshift=2ex]\tikztostart.east)
        |- (X) [near end]\tikztonodes
        -| ([xshift=-2ex]\tikztotarget.west)
        -- (\tikztotarget)}] \\
      & \cohomology{i}(X; \sheaf{G}) \arrow[r] & \cohomology{i}(V; \sheaf{G}|_{V}) \oplus \cohomology{i}(W; \sheaf{G}|_{W}) \arrow[r] \arrow[d, phantom, ""{coordinate, name=Y}] &  \cohomology{i}(V \cap W; \sheaf{G}|_{V \cap W}) \arrow[dll, rounded corners, to path={
        -- ([xshift=2ex]\tikztostart.east)
        |- (Y) [near end]\tikztonodes
        -| ([xshift=-2ex]\tikztotarget.west)
        -- (\tikztotarget)}] \\
      & \cohomology{i+1}(X; \sheaf{G}) \arrow[r] & \cdots
    \end{tikzcd}
  \end{center}

  Note that the \hyperref[def:stalk]{stalks} of \(\sheaf{G}\) are
  \begin{equation*}
    \sheaf{G}_{p} = \begin{cases}
      \sheaf{F}_{p} & p \in U \\
      0 & p \notin U
    \end{cases}
  \end{equation*}
  Therefore since \(W \subseteq A\) and \(A \cap U = \varnothing\) we have \(\sheaf{G}_{p} = 0\) for all \(p \in W\).
  In particular \(\sheaf{G}|_{W} = 0\) and thus \(\sheaf{G}|_{V \cap W} = 0\).
  Plugging this into the above \hyperref[def:exact-sequence]{long exact sequence} we get \hyperref[def:categorical-isomorphism]{isomorphisms}
  \begin{equation*}
    \cohomology{i}(X; \sheaf{G}) \isoarrow \cohomology{i}(V; \sheaf{G}|_{V})
  \end{equation*}
  To finish, we need to check that \(\cohomology{i}(V; \sheaf{G}) = \cohomology{i} \left( X \setminus Z, A \setminus Z, \sheaf{F} \right)\).
  Note that
  \begin{equation*}
    V \setminus ( A \setminus Z ) = \left( X \setminus Z \right) \setminus \left( A \setminus Z \right) = X \setminus A = U
  \end{equation*}
  Consider the open immersion \(j'\colon U \hookrightarrow V\).
  Since \(\sheaf{G}\) was formed by extending \(\sheaf{F}|_{U}\) by zero to \(X\), restricting it to \(V\) (which fully contains \(U\)) yields the same extension by zero relative to \(V\):
  \begin{equation*}
    \sheaf{G}|_{V} = j'_{!} \left( \sheaf{F}|_{U} \right) = j'_{!} j'^{-1} \left( \sheaf{F}|_{V} \right)
  \end{equation*}
  Therefore
  \begin{equation*}
    \cohomology{i} \left( V; \sheaf{G}|_{V} \right) = \cohomology{i} \left( V; j'_{!}j'^{-1}\sheaf{F}|_{V} \right) = \cohomology{i} \left( X \setminus Z, A \setminus Z; \sheaf{F} \right) \qedhere
  \end{equation*}
\end{proof}

\begin{theorem}[Exactness]\label{thm:exactness-axiom}
  Let \(A \subseteq X\) be a closed subset and \(i\colon A \hookrightarrow X\) the closed immersion.
  For any \hyperref[def:sheaf]{sheaf of \(\mathbf{k}\)-modules} \(\sheaf{F}\), there exists a \hyperref[def:exact-sequence]{long exact sequence} in cohomology:
  \begin{center}
    \begin{tikzcd}[column sep=small]
      \cdots \arrow[r] & \cohomology{i}(X, A; \sheaf{F}) \arrow[r] & \cohomology{i}(X; \sheaf{F}) \arrow[r] & \cohomology{i}(A; \sheaf{F}_{A}) \arrow[r] & \cohomology{i+1}(X, A; \sheaf{F}) \arrow[r] & \cdots
    \end{tikzcd}
  \end{center}
\end{theorem}
\begin{proof}
  By \autoref{thm:SES-exactness} we have a \hyperref[def:exact-sequence]{short exact sequence} of \hyperref[def:sheaf]{sheaves}:
  \begin{equation*}
    0 \to j_{!}j^{-1}\sheaf{F} \to \sheaf{F} \to i_{\ast}i^{-1}\sheaf{F} \to 0
  \end{equation*}
  Applying the \hyperref[thm:LES-derived-functor]{long exact sequence of derived functors} to \(\Gamma(X, -)\) and using \autoref{prop:cohomology-subsets} for the term \(\cohomology{i}(X; i_{\ast}i^{-1}\sheaf{F}) \cong \cohomology{i}(A; \sheaf{F}_{A})\), we obtain the result.
\end{proof}

\begin{theorem}[Additivity]\label{thm:additivity-axiom}
  Let \(X = \coprod_{\alpha \in \Lambda} X_{\alpha}\) be a disjoint union of topological spaces.
  Then for any \hyperref[def:sheaf]{sheaf of \(\mathbf{k}\)-modules} \(\sheaf{F}\) on \(X\):
  \begin{equation*}
    \cohomology{i}(X; \sheaf{F}) \cong \prod_{\alpha \in \Lambda} \cohomology{i}(X_{\alpha}; \sheaf{F}|_{X_{\alpha}})
  \end{equation*}
\end{theorem}
\begin{proof}
  By definition of the disjoint union topology, each \(X_{\alpha}\) is an open subset of \(X\).

  Let \(\sheaf{F} \in \Mod ( \underline{\mathbf{k}}_{X} )\).
  Since \(X_{\alpha} \cap X_{\beta} = \varnothing\) if \(\alpha \neq \beta\) we have
  \begin{equation*}
    \Gamma(X, \sheaf{F}) \cong \prod_{\alpha \in \Lambda} \Gamma(X_{\alpha}, \sheaf{F}|_{X_{\alpha}})
  \end{equation*}

  Furthermore if \(\sheaf{F} \to \sheaf{I}^{\bullet}\) is an \hyperref[def:resolution]{injective resolution} since \(j_{\alpha}^{-1}\) is \hyperref[def:adjoint-functors]{right adjoint} to \(j_{\alpha, !}\), \hyperref[cor:extension-by-zero-sheaf-exact]{which is exact}, by \autoref{prop:adjoints-and-injectives} \(\sheaf{I}^{\bullet}|_{X_{\alpha}}\) is an \hyperref[def:resolution]{injective resolution} of \(\sheaf{F}|_{X_{\alpha}}\).
  
  Therefore
  \begin{align*}
    \cohomology{i}(X; \sheaf{F}) &\cong \operatorname{R}^{i}\Gamma(X, \sheaf{I}^{\bullet}) \\
                                 &\cong  \operatorname{R}^{i} \left( \prod_{\alpha \in \Lambda} \Gamma(X_{\alpha}, \sheaf{I}^{\bullet}|_{X_{\alpha}}) \right) \\
                                 &\cong \prod_{\alpha \in \Lambda} \operatorname{R}^{i} \Gamma(X_{\alpha}, \sheaf{I}^{\bullet}|_{X_{\alpha}}) \\
                                 &\cong \prod_{\alpha \in \Lambda} \cohomology{i}(X_{\alpha}, \sheaf{F}|_{X_{\alpha}}) \qedhere
  \end{align*}
\end{proof}

% \begin{theorem}[Dimension]
%   Let \(\{\ast\}\) be the one point topological space.
%   For all \hyperref[def:sheaf]{sheaves of \(\mathbf{k}\)-modules} \(\sheaf{F} \in \Mod(\underline{\mathbf{k}}_{X})\) and \(n \neq 0\)
%   \begin{equation*}
%     \cohomology{n}(\{\ast\}; \sheaf{F}) \cong 0
%   \end{equation*}
% \end{theorem}
% \begin{proof}
%   We have
%   \begin{align*}
%     \cohomology{n}( \{\ast\}; \sheaf{F} ) = \operatorname{R}^{n} \Gamma(\{\ast\}, \sheaf{F})
%   \end{align*}
%   Note that
%   \begin{equation*}
%     \Gamma(\{\ast\}, \sheaf{F}) = \sheaf{F}(\{\ast\}) = \sheaf{F}_{\ast}
%   \end{equation*}
%   Therefore since the stalk functor is exact for all \(n \neq 0\)
%   \begin{equation*}
%     \cohomology{n} (\{\ast\}; \sheaf{F}) = \operatorname{R}^{n} \sheaf{F}_{\ast} = 0 \qedhere
%   \end{equation*}
% \end{proof}

Having fully verified all five axioms by the \hyperref[thm:uniqueness-cohomology-theories]{uniqueness of cohomology theories} we obtain complete equivalences between \hyperref[def:sheaf-cohomology]{sheaf cohomology} and classical cohomology groups.

In particular from \autoref{prop:manifold-is-CW-complex} and the \hyperref[thm:homotopy-invariance-cohomology]{homotopy invariance of cohomology} we have the following equivalence.
\begin{corollary}
  Let \(X\) be a space homotopic to a \hyperref[def:CW-complex]{CW-complex}.
  For any \hyperref[def:constant-sheaf]{constant sheaf of abelian groups} \(\underline{G} \in \Mod(\underline{\mathbb{Z}}_{X})\) we have an \hyperref[def:categorical-isomorphism]{isomorphism}
    \begin{equation*}
        \cohomology{i}(X; \underline{G}) \cong \operatorname{H}_{\operatorname{sing}}^{i}(X; G)
    \end{equation*}
\end{corollary}

Furthermore, we can combine this with the work we did in the last chapter obtaining a proof that \hyperref[def:de-rham-cohomology]{de Rham cohomology} and singular cohomology coincide for compact smooth manifolds.
\begin{corollary}
  Let \(X\) be a compact smooth manifold.
  We have \hyperref[def:categorical-isomorphism]{isomorphisms} (of abelian groups)
  \begin{equation*}
    \operatorname{H}_{\operatorname{sing}}^{i}(X; \mathbb{R}) \cong \cohomology{i}(X; \underline{\mathbb{R}}) \cong \operatorname{H}_{\operatorname{dR}}(X; \mathbb{R})
  \end{equation*}
\end{corollary}

\end{document}
